
\documentclass{article}
\usepackage{colortbl}
\usepackage{makecell}
\usepackage{multirow}
\usepackage{supertabular}

\begin{document}

\newcounter{utterance}

\centering \large Interaction Transcript for game `cladder', experiment `full\_v1.5\_default', episode 6687 with score{-}Qwen3.5{-}9B{-}sp.
\vspace{24pt}

{ \footnotesize  \setcounter{utterance}{1}
\setlength{\tabcolsep}{0pt}
\begin{supertabular}{c@{$\;$}|p{.15\linewidth}@{}p{.15\linewidth}p{.15\linewidth}p{.15\linewidth}p{.15\linewidth}p{.15\linewidth}}
    \# & \multicolumn{2}{c}{Player} && \multicolumn{2}{c}{Game Master} \\
    \hline

    \theutterance \stepcounter{utterance}  
    & & & \multicolumn{4}{p{0.6\linewidth}}{
        \cellcolor[rgb]{0.9,0.9,0.9}{
            \makecell[{{p{\linewidth}}}]{
                \texttt{\tiny{[P1$\langle$GM]}}
                \texttt{You are an expert in causal inference. The following question is not a typical commonsense query, but rather a meticulously designed question created by a professor specializing in causal inference, intended to assess the students' mastery of the course content.} \\
\texttt{Imagine a self{-}contained, hypothetical world with only the following conditions, and without any unmentioned factors or causal relationships: Xevo has a direct effect on tijw and tijv. Tijv has a direct effect on gyzp. Tijw has a direct effect on gyzp. The overall probability of xevo is 48\%. For those who are not xevo, the probability of gyzp is 39\%. For those who are xevo, the probability of gyzp is 78\%. Is gyzp more likely than not gyzp overall?} \\
\texttt{Start your answer with 'Yes' or 'No', followed by additional reasoning or evidence to support your explanation.} \\
            }
        }
    }
     \\ \\

    \theutterance \stepcounter{utterance}  
    & \multicolumn{4}{p{0.6\linewidth}}{
        \cellcolor[rgb]{0.561,0.737,0.561}{
            \makecell[{{p{\linewidth}}}]{
                \texttt{\tiny{[P1$\rangle$GM]}}
                \texttt{Yes} \\
\\ 
\texttt{**Reasoning:**} \\
\\ 
\texttt{To determine if \$gyzp\$ is more likely than not \$gyzp\$ overall, we must calculate the marginal probability \$P(gyzp)\$ using the Law of Total Probability. This law states that the total probability of an event is the sum of the probabilities of that event occurring under each mutually exclusive condition, weighted by the probability of those conditions.} \\
\\ 
\texttt{Here, the conditions are whether \$xevo\$ occurs or not. We are given:} \\
\texttt{1.  \$P(xevo) = 0.48\$} \\
\texttt{2.  Therefore, \$P(\textbackslash{}neg xevo) = 1 {-} 0.48 = 0.52\$} \\
\texttt{3.  \$P(gyzp | xevo) = 0.78\$} \\
\texttt{4.  \$P(gyzp | \textbackslash{}neg xevo) = 0.39\$} \\
\\ 
\texttt{The causal structure described (Xevo \$\textbackslash{}to\$ tijw, Xevo \$\textbackslash{}to\$ tijv, tijw \$\textbackslash{}to\$ gyzp, tijv \$\textbackslash{}to\$ gyzp) explains *why* the probabilities differ based on the state of \$xevo\$, but for the specific question of the overall likelihood, we rely on the provided conditional probabilities. The direct effects imply that \$xevo\$ is a common cause influencing \$gyzp\$ through mediators, which is consistent with the provided conditional data.} \\
\\ 
\texttt{Calculating the overall probability:} \\
\texttt{\$\$P(g} \\
            }
        }
    }
    & & \\ \\

    \theutterance \stepcounter{utterance}  
    & & & \multicolumn{2}{p{0.3\linewidth}}{
        \cellcolor[rgb]{0.9,0.9,0.9}{
            \makecell[{{p{\linewidth}}}]{
                \texttt{\tiny{[GM$|$GM]}}
                \texttt{Yes} \\
            }
        }
    }
    & & \\ \\

    \theutterance \stepcounter{utterance}  
    & & & \multicolumn{2}{p{0.3\linewidth}}{
        \cellcolor[rgb]{0.9,0.9,0.9}{
            \makecell[{{p{\linewidth}}}]{
                \texttt{\tiny{[GM$|$GM]}}
                \texttt{yes} \\
            }
        }
    }
    & & \\ \\

    \theutterance \stepcounter{utterance}  
    & & & \multicolumn{2}{p{0.3\linewidth}}{
        \cellcolor[rgb]{0.9,0.9,0.9}{
            \makecell[{{p{\linewidth}}}]{
                \texttt{\tiny{[GM$|$GM]}}
                \texttt{game\_result = WIN} \\
            }
        }
    }
    & & \\ \\

\end{supertabular}
}

\end{document}
